\section*{摘要}

去中心化金融（DeFi）自2020年以来经历了指数级增长，以Aave、Compound和MakerDAO为代表的借贷协议在多个区块链网络上管理的总锁仓价值峰值超过1000亿美元。这些协议完全通过智能合约运行，自动管理抵押品、借贷和清算，从而产生所有参与者行为的完全透明链上记录。然而，现有关于DeFi借贷的学术文献主要关注协议层面的风险暴露、清算机制和整体市场效率，对微观层面个体借款人在其头寸接近风险阈值时的行为响应关注较少。本研究报告普查了三个相互关联领域的相关文献——DeFi借贷基础设施、行为金融与前景理论以及链上信用风险评估——并识别出三者交叉处一个明确的、尚未充分探索的研究空白。我们提出了一项研究，探讨借款人在头寸风险上升期间是否采取了可提供超越传统链上指标（如健康因子、抵押率和账户活动）的风险信息的主动调整行为。所提出的方法论从公开可用的区块链数据中重构了头寸风险轨迹，区分了借款人主动操作与被动清算事件，并利用判别框架检验了行为过程变量的增量预测贡献。本文讨论了可能的成果情景，包括对DeFi协议设计和信用风险建模的贡献，并对研究的边界条件进行了批判性评估。

\vspace{8pt}
\noindent\textbf{关键词：} 去中心化金融，DeFi借贷，行为金融，前景理论，信用风险，区块链分析，链上数据
